% =============================================================================
%  Partie I — La théorie logique (Chapitre I)
% =============================================================================
\section{Le langage et les schémas logiques (I.1--I.4)}

La base de la tour est la \emph{théorie logique} : un langage formel et sept
schémas d'axiomes, sans aucun axiome propre aux ensembles. Tout ce qui s'y
démontre vaut dans \emph{toute} théorie plus forte. C'est ici, et nulle part
ailleurs, qu'on prouve les lois du calcul propositionnel et des quantificateurs.

\subsection{Le langage : assemblages, $\tau$, signes (I.1)}

Un objet est un \emph{assemblage de signes}, soit un \emph{terme} (il désigne un
objet), soit une \emph{relation} (elle est vraie ou fausse). Les signes
primitifs sont logiques ($\neg$, $\vee$, et l'opérateur $\tau$), relationnels
($=$, $\in$), et des lettres. Bourbaki ne prend pas les quantificateurs comme
primitifs : il prend le \textbf{$\tau$ de Hilbert}. Pour une relation $R$, le
terme $\tau_x(R)$ désigne « un objet $x$ tel que $R$, s'il en existe un ». Tout
le reste en est \emph{défini} :
\[
  (\exists x)\,R \;:=\; \bigl(\tau_x(R)\mid x\bigr)R,
  \qquad
  (\forall x)\,R \;:=\; \neg(\exists x)\,\neg R,
\]
\[
  R\Rightarrow S \;:=\; (\neg R)\vee S,
  \quad
  R\wedge S \;:=\; \neg(\neg R\vee\neg S),
  \quad
  R\Leftrightarrow S \;:=\; (R\Rightarrow S)\wedge(S\Rightarrow R).
\]
Le module \texttt{assemblage} réalise les assemblages comme structures
immuables et hachables (égalité et hachage en $O(1)$ après interning), la
substitution $(T\mid x)R$, et les prédicats de lecture
($\texttt{est\_terme}$, $\texttt{est\_relation}$) qui valident la bonne
formation. Conséquence notable : le $\tau$ encode l'\emph{axiome du choix} dans
le langage même --- la théorie est « choisie » dès le départ.

\subsection{Les sept schémas (I.3, I.5)}

Les seuls axiomes logiques sont les schémas \texttt{s1}--\texttt{s7} du noyau
(PDF p.~33--38) :
\[
\begin{array}{ll}
  \text{S1} & (R\vee R)\Rightarrow R \\
  \text{S2} & R\Rightarrow(R\vee S) \\
  \text{S3} & (R\vee S)\Rightarrow(S\vee R) \\
  \text{S4} & (R\Rightarrow S)\Rightarrow\bigl((T\vee R)\Rightarrow(T\vee S)\bigr) \\
  \text{S5} & (T\mid x)R \Rightarrow (\exists x)R \\
  \text{S6} & (T=U)\Rightarrow\bigl((T\mid x)R \Leftrightarrow (U\mid x)R\bigr) \\
  \text{S7} & (\forall x)(R\Leftrightarrow S)\Rightarrow \tau_x(R)=\tau_x(S) \\
\end{array}
\]
S1--S4 sont le calcul propositionnel ; S5 gouverne le $\tau$ (un témoin suffit
à l'existence) ; S6--S7 sont les schémas \emph{égalitaires} (substitutivité de
Leibniz et extensionnalité du $\tau$). Dans le noyau, chacun est une fonction
qui \emph{construit} l'assemblage du schéma et le scelle en \theoreme{} clos ---
aucune théorie en argument : ce sont les briques universelles.

\subsection{La règle d'inférence et les critères dérivés (I.2)}

Le noyau ne possède qu'\emph{une} règle d'inférence de base, le
\textbf{détachement} (critère C1) :
\[
  \Gamma\vdash R \quad\text{et}\quad \Delta\vdash(R\Rightarrow S)
  \qquad\Longrightarrow\qquad \Gamma\cup\Delta\vdash S.
\]
Le \texttt{modus\_ponens} du noyau le réalise \emph{honnêtement} : il décode la
majeure comme l'assemblage $\vee\,\neg R\,S$, vérifie que la mineure coïncide
avec $R$, puis reconstruit $S$ et réunit les hypothèses. Deux autres règles sont
admises comme \emph{primitives de confiance}, car leur validité est un critère
que Bourbaki démontre : la \textbf{déduction} $\texttt{loi\_deduction}$ (C6 :
de $\Gamma\vdash B$ déduire $\Gamma\setminus\{A\}\vdash(A\Rightarrow B)$) et la
\textbf{généralisation} $\texttt{generalisation}$ (C27 : de $\Gamma\vdash R$,
avec $x$ non libre dans $\Gamma$, déduire $\Gamma\vdash(\forall x)R$). Tout le
reste des critères C de Bourbaki (raisonnement par l'absurde, contraposition,
transitivité de $\Rightarrow$, échanges de quantificateurs) est \emph{dérivé} au
sein de la théorie logique, par composition de ces primitives --- jamais ajouté
à la base de confiance.

\subsection{Ce qui est formalisé dans la théorie logique}

Vivent dans cette couche, et y sont prouvés \emph{sans aucun axiome ensembliste} :
les lois propositionnelles usuelles (idempotence, commutativité/associativité
de $\vee$ et $\wedge$, lois de De~Morgan logiques, tiers exclu, double
négation), les règles de la déduction et de la généralisation comme lemmes
réutilisables, et les manipulations de quantificateurs (introduction/élimination
du $\forall$ et du $\exists$, échange de quantificateurs de même nature). Le
lemme propositionnel $\neg(P\Leftrightarrow\neg P)$, par exemple, est construit
ici --- et c'est lui qui, une fois remonté en théorie des ensembles, interdit
l'ensemble de Russell (Partie~II). C'est l'illustration de l'héritage vers le
haut : un fait logique, prouvé une seule fois au plus bas, sert tel quel à tous
les étages supérieurs.

\subsection{Quantificateurs typiques : le critère C41 (I.4.4)}
Les quantificateurs \emph{typiques} (relativisés à une relation $A$) sont
$(\exists_A x)R:=(\exists x)(A\wedge R)$ et $(\forall_A x)R:=\neg(\exists x)(A\wedge\neg R)$.
Leurs critères $C39$ (monotonie), $C40$ (distribution sur $\wedge/\vee$) et $C42$
(commutation) sont formalisés clos ; le présent jalon ajoute le \textbf{critère
$C41$} (E~I.37), seul absent de la série, pour une lettre $x$ \emph{ne figurant pas
dans $R$} :
\[
  (\forall_A x)(R\vee S)\Leftrightarrow\bigl(R\vee(\forall_A x)S\bigr),\qquad
  (\exists_A x)(R\wedge S)\Leftrightarrow\bigl(R\wedge(\exists_A x)S\bigr).
\]
\textbf{Preuve (\Nstar).} Le cas $\exists$ est l'extraction du facteur $R$ via le
critère $C33$ (\texttt{et\_existe\_droite}, $x\notin R$) précédé d'un réarrangement
de la conjonction sous $(\exists x)$ ; le cas $\forall$ en est le \emph{dual} par
De~Morgan (sur $\neg R,\neg S$) et double négation. Clos (théorèmes purs), la
fraîcheur $x\notin R$ étant load-bearing (refus explicite sinon). Module
\texttt{.../i\_3\_quantifies/criteres\_typiques\_c41.py}.
